\SourcePage{274}{277}
\section{Time reversal and Kramers' theorem}

In quantum mechanics, symmetry of motion under reversal of the sign of time
means the following: if $\psi$ is the wave function of some stationary state
of a system, then the ``time-reversed'' wave function, denoted by
$\psi^{\text{rev}}$, also describes a possible state with
\SourcePage{275}{278}
the same energy. It was stated at the end of \secref{18} that
$\psi^{\text{rev}}$ coincides with the complex-conjugate function $\psi^*$.
In this simple form, the statement applies to wave functions when particle
spin is disregarded. With spin present, it requires qualification.

Represent the wave function of a spin-$s$ particle by the contravariant
spinor $\psi^{\lambda\mu\ldots}$ of rank $2s$. On passing to the
complex-conjugate functions $\psi^{\lambda\mu\ldots*}$, however, we obtain
quantities that transform as components of a covariant spinor. Time reversal
therefore replaces $\psi^{\lambda\mu\ldots}$ by a new wave function whose
covariant components are defined by
\begin{equation}
 \psi^{\text{rev}}_{\lambda\mu\ldots}=\psi_{\lambda\mu\ldots}^*.
 \tag{60.1}
\end{equation}
For a specified set of index values $\lambda,\mu,\ldots$, the covariant and
contravariant spinor components correspond to angular-momentum projections of
opposite sign. In terms of the functions $\psi_{s\sigma}$, time reversal thus
replaces $\psi_{s\sigma}$ by $\psi_{s,-\sigma}$, as it must, since reversing
time reverses the direction of angular momentum. The precise correspondence
following from (60.1) is
\begin{equation}
 \psi^{\text{rev}}_{s,-\sigma}=\psi_{s\sigma}^*(-1)^{s-\sigma}.
 \tag{60.2}
\end{equation}
In other words, the replacement
$\psi_{s\sigma}\longrightarrow\psi_{s\sigma}^*$ required by time reversal
means the replacement\footnote{Notice the agreement of the complex-conjugation
rule for a spherical function, (28.9), with the general rule (60.3).}
\begin{equation}
 \psi_{s\sigma}\longrightarrow\psi_{s,-\sigma}(-1)^{s-\sigma}.
 \tag{60.3}
\end{equation}
Applying this operation twice gives
\[
 \psi_{s\sigma}\longrightarrow\psi_{s,-\sigma}(-1)^{s-\sigma}
 \longrightarrow\psi_{s\sigma}(-1)^{s-\sigma}(-1)^{s+\sigma}
 =\psi_{s\sigma}(-1)^{2s}.
\]
Thus two successive time reversals restore the original wave function only
for integral spin; for half-integral spin they change its sign.

Consider an arbitrary system of interacting particles. When relativistic
interactions are included, the orbital and spin angular momenta of such a
system are not in general conserved separately; only the total angular
momentum $J$ is conserved. In the absence of an external field, each energy
level of the system is $(2j+1)$-fold degenerate. An external field generally
removes this degeneracy. The question is whether the degeneracy can be removed
completely, leaving the system with only nondegenerate levels. This question
is closely connected with time-reversal symmetry.
\SourcePage{276}{279}
Classical electrodynamics is invariant under reversal of the sign of time if
the electric field is left unchanged and the sign of the magnetic field is
reversed.\footnote{See II, \S~17; see also the remark at the end of \secref{111}.}
This fundamental property of motion must persist in quantum mechanics.
Time-reversal symmetry therefore holds not only for a closed system, but also
in any external electric field, provided no magnetic field is present.

The wave functions of the system are spinors $\psi^{\lambda\mu\ldots}$ of
rank $n$ equal to twice the sum of the spins of all particles,
$n=2\sum s_a$; this sum need not coincide with the total spin $S$ of the
system. From the foregoing, in an arbitrary electric field a wave function
and its time reverse must correspond to states of equal energy. A necessary
condition for a level to be nondegenerate is that these states be identical,
so their wave functions must agree apart from a constant factor. Both must,
of course, be expressed as spinors of the same kind, covariant or
contravariant.

Write $\psi^{\text{rev}}_{\lambda\mu\ldots}
=C\psi_{\lambda\mu\ldots}$, or, by (60.1),
\begin{equation}
 \psi_{\lambda\mu\ldots}^*=C\psi_{\lambda\mu\ldots},
 \tag{60.4}
\end{equation}
where $C$ is constant. Complex conjugation of the two sides gives
\[
 \psi^{\lambda\mu\ldots}=C^*\psi^{\lambda\mu\ldots*}.
\]
Lower the indices on the left and correspondingly raise those on the right.
This means multiplying both sides by
$g_{\alpha\lambda}g_{\beta\mu}\ldots$ and summing over
$\lambda,\mu,\ldots$; on the right one must use
\[
 g_{\alpha\lambda}g_{\beta\mu}\ldots
 =(-1)^n g_{\lambda\alpha}g_{\mu\beta}\ldots.
\]
The result is
\[
 \psi_{\lambda\mu\ldots}=C^*(-1)^n\psi_{\lambda\mu\ldots}^*.
\]
Substituting $\psi_{\lambda\mu\ldots}^*$ from (60.4), we find
\[
 \psi_{\lambda\mu\ldots}=(-1)^nCC^*\psi_{\lambda\mu\ldots}.
\]
\SourcePage{277}{280}
This equality must hold identically, so $(-1)^nCC^*=1$. But $|C|^2$ is
necessarily positive; the condition is therefore possible only for even $n$,
that is, for an integral value of $\sum s_a$. For odd $n$, or half-integral
$\sum s_a$,\footnote{When $\sum s_a$ is integral (half-integral), all possible
values of the total spin $S$ of the system are likewise integral
(half-integral).} condition (60.4) cannot be satisfied.

We thus reach the result that an electric field can remove the degeneracy
completely only in a system with an integral sum of particle spins. In a
system with a half-integral spin sum, every level in an arbitrary electric
field must be doubly degenerate; the two distinct states of equal energy
correspond to complex-conjugate spinors.\footnote{If the electric field has
high (cubic) symmetry, fourfold degeneracy may also occur (see \secref{99} and its
problem).} (H.~A. Kramers, 1930).

We add one mathematical remark. A relation of the form (60.4), with real
$C$, is mathematically the condition that some set of real quantities can be
put in correspondence with the spinor components; it may be called a
``reality'' condition for the spinor.\footnote{It is meaningless to speak of
a spinor as literally real, since complex-conjugate spinors have different
transformation laws.} The impossibility of (60.4) for odd $n$ means that no
real quantity can correspond to a spinor of odd rank. For even $n$, by
contrast, (60.4) can hold, and $C$ can be real. In particular, a real vector
can be put in correspondence with a symmetric second-rank spinor when (60.4)
holds with $C=1$:
\[
 \psi^{\lambda\mu*}=\psi_{\lambda\mu}
\]
(as is easily verified from (57.8), (57.9)). In general, (60.4) with $C=1$
is the ``reality'' condition for a symmetric spinor of any even rank.
